\documentclass[12pt,a4paper]{article}
\usepackage[portuguese]{babel}
\usepackage[utf8]{inputenc}
\usepackage[T1]{fontenc}
\usepackage{graphicx}
\usepackage{amsmath}
\usepackage{amssymb}
\usepackage{listings}
\usepackage{xcolor}
\usepackage{hyperref}
\usepackage{geometry}
\usepackage{multirow}
\geometry{margin=2.5cm}


% Configuração para código C
\lstset{
    language=C,
    basicstyle=\ttfamily\small,
    keywordstyle=\color{blue},
    commentstyle=\color{gray},
    stringstyle=\color{red},
    numbers=left,
    numberstyle=\tiny\color{gray},
    stepnumber=1,
    numbersep=5pt,
    backgroundcolor=\color{white},
    showspaces=false,
    showstringspaces=false,
    showtabs=false,
    frame=single,
    tabsize=2,
    captionpos=b,
    breaklines=true,
    breakatwhitespace=false
}

\title{Relatório do Projeto AED}
\author{Tiago Pedrosa : 116113 \\ Gustavo Rodrigues : 126033 }
\date{Algoritmos e Estruturas de Dados \\ Universidade de Aveiro \\ Engenharia de Computadores e Informática}

\begin{document}

\maketitle
\newpage

\tableofcontents
\newpage

\section{Introdução}

Este relatório apresenta o desenvolvimento do trabalho prático de Algoritmos e Estruturas de Dados, que consiste na implementação de um módulo para processamento de imagens RGB utilizando uma tabela de cores (LUT).

O projeto utiliza uma representação eficiente de imagens RGB onde cada pixel armazena um índice para uma tabela de cores. Esta abordagem é vantajosa quando o número de cores distintas na imagem é significativamente menor que o número total de pixels.

\subsection{Objetivos}

Este trabalho tem como objetivos implementar funções de manipulação de imagens e desenvolver algoritmos de preenchimento de regiões utilizando abordagens diferentes. Inclui também a implementação de segmentação de imagens baseada em regiões conectadas e análise da complexidade computacional das funções desenvolvidas.

\section{Estruturas de Dados e Funções Básicas}

\subsection{Estrutura da Imagem}

A estrutura da imagem RGB com LUT é definida da seguinte forma:

\begin{lstlisting}
struct image {
    uint32 width;        // Largura da imagem
    uint32 height;       // Altura da imagem
    uint16** image;      // Matriz de labels (indices para LUT)
    uint16 num_colors;   // Numero de cores utilizadas
    rgb_t* LUT;          // Look-Up Table (tabela de cores)
};
\end{lstlisting}

\subsection{Funções Implementadas}

\subsubsection{ImageCopy}

A função \texttt{ImageCopy} cria uma cópia de uma imagem. O algoritmo aloca uma nova estrutura com as mesmas dimensões, copia a LUT completa e depois copia todos os pixels linha por linha. A complexidade temporal é $O(W \times H + C)$, onde $W$ é a largura, $H$ é a altura e $C$ é o número de cores.

\subsubsection{ImageRotate90 e ImageRotate180}

A rotação de 90° transforma um pixel em $(u, v)$ para $(H-1-v, u)$ na nova imagem, invertendo as dimensões. A rotação de 180° transforma um pixel em $(u, v)$ para $(W-1-u, H-1-v)$ mantendo as dimensões. Ambas apresentam complexidade $O(W \times H)$.

\section{Análise Detalhada da Função ImageIsEqual}

\subsection{Análise Formal da Complexidade}

A função \texttt{ImageIsEqual} verifica se duas imagens são iguais comparando as cores RGB de todos os pixels. O algoritmo verifica inicialmente se as dimensões são iguais; posteriormente compara pixel a pixel até encontrar uma diferença ou terminar a varredura.

Para duas imagens de dimensões $W \times H$, a análise de complexidade é:

\begin{itemize}
\item \textbf{Melhor caso:} $\Theta(1)$ - ocorre quando as dimensões das imagens são diferentes, terminando imediatamente sem comparação de pixels.

\item \textbf{Pior caso:} $\Theta(W \times H)$ - ocorre quando as imagens são completamente iguais, sendo necessário comparar todos os $W \times H$ pixels até ao fim. Cada comparação envolve 2 acessos aos arrays de pixels e comparação das cores RGB.

\item \textbf{Caso médio:} $O(W \times H)$ - assumindo imagens com dimensões iguais mas conteúdo potencialmente diferente. Se as imagens diferem, a função termina quando encontra a primeira diferença. Estatisticamente, espera-se que a diferença ocorra em média após verificar metade dos pixels, mas em notação assintótica o fator constante é irrelevante.
\end{itemize}

Em termos de acessos à memória (PIXMEM), o pior caso requer $2 \times W \times H$ acessos aos pixels (uma leitura em cada imagem por comparação). A otimização early exit permite reduzir significativamente este número quando diferenças são encontradas cedo.

\subsection{Análise Experimental}

Realizaram-se testes computacionais com imagens de diferentes tamanhos e conteúdos, registando o número de acessos PIXMEM:

\begin{table}[h]
\centering
\begin{tabular}{|l|c|c|c|c|}
\hline
\textbf{Teste} & \textbf{Dimensão} & \textbf{Conteúdo} & \textbf{PIXMEM} & \textbf{Fórmula} \\
\hline
1 & $100 \times 100$ & Xadrez (iguais) & 20,000 & $2 \times 100 \times 100$ \\
\hline

2 & $24 \times 7$ & feep.pbm (iguais) & 336 & $2 \times 24 \times 7$ \\
\hline
3 & $50 \times 50$ & Branca (iguais) & 5,000 & $2 \times 50 \times 50$ \\
\hline
4 & $40 \times 60$ & Xadrez (iguais) & 4,800 & $2 \times 40 \times 60$ \\
\hline
\end{tabular}
\caption{Resultados experimentais da função ImageIsEqual}
\end{table}

\subsection{Comparação: Análise Experimental vs Formal}

Os resultados experimentais confirmam perfeitamente a análise formal. Para imagens iguais, o número de acessos PIXMEM corresponde exatamente à fórmula $2 \times W \times H$, validando a complexidade $\Theta(W \times H)$ do pior caso.
\newpage
\section{Algoritmos de Preenchimento de Regiões}

\subsection{Descrição dos Algoritmos}

Implementaram-se três versões do algoritmo flood-fill para preenchimento de regiões conexas com conectividade de 4 vizinhos:

\subsubsection{Versão Recursiva (ImageRegionFillingRecursive)}

O algoritmo verifica se o pixel é válido e tem a cor original. Para cada pixel válido, pinta o pixel com a nova cor e chama recursivamente a função para os 4 vizinhos (cima, baixo, esquerda e direita).

\subsubsection{Versão com Stack Explícito (ImageRegionFillingWithSTACK)}

O algoritmo cria um stack e empilha o pixel inicial. Enquanto o stack não estiver vazio, desempilha um pixel, verifica se é válido e tem a cor original, pinta o pixel com a nova cor e empilha os 4 vizinhos.

\subsubsection{Versão com Fila (ImageRegionFillingWithQUEUE)}

O algoritmo cria uma fila e enfileira o pixel inicial. Enquanto a fila não estiver vazia, retira um pixel, verifica se é válido e tem a cor original, pinta o pixel com a nova cor e enfileira os 4 vizinhos. 

\subsection{Análise de Complexidade}

\begin{table}[h]
\centering
\begin{tabular}{|l|c|c|c|c|}
\hline
\textbf{Algoritmo} & \textbf{Tempo} & \textbf{Espaço} & \textbf{Ordem de Preenchimento} \\
\hline
Recursivo & $O(N)$ & $O(N)$  & Profundidade (stack do sistema) \\
STACK & $O(N)$ & $O(N)$  & Profundidade (stack explícito) \\
QUEUE & $O(N)$ & $O(N)$ & Largura (níveis sucessivos)

 \\
\hline
\end{tabular}
\caption{Comparação dos algoritmos de flood-fill}
\end{table}

Todos os algoritmos visitam exatamente $N$ pixels (onde $N$ é o número de pixels na região) e realizam $O(1)$ operações por pixel, resultando em $O(N)$ operações totais. A complexidade espacial é $O(N)$ em todos os casos: na versão recursiva, a profundidade máxima da pilha de recursão é $O(N)$; nas versões iterativas, a estrutura de dados (stack ou fila) armazena no máximo $O(N)$ elementos.

\subsection{Comparação das Estratégias}

\subsubsection{Versão Recursiva}

 Vantagens: simplicidade e clareza. Desvantagens: risco de stack overflow para regiões grandes (a pilha de sistema tem espaço limitado) e overhead de chamadas de função recursiva.
\subsubsection{Versão com Stack}

 Vantagens: elimina o risco de stack overflow do sistema, mesmo para regiões grandes; permite alocação dinâmica de memória. Desvantagens: mais código, overhead de gerenciamento de stack. 

\subsubsection{Versão com Fila}

 Vantagens: preenchimento em camadas sucessivas a partir do ponto inicial, útil para análise de propagação. Desvantagens: maior consumo de memória em regiões muito alongadas. 

\subsection{Resultados Experimentais}
\begin{table}[h]
\begin{tabular}{|l|c|c|c|c|c|}
\hline
\textbf{Teste} & \textbf{Algoritmo} & \textbf{Pixels Preench.} & \textbf{PIXMEM} & \textbf{Tempo ($\mu$s)}\\
\hline
\multirow{3}{*}{Imagem $50 \times 50$ branca} & Recursivo & 2.500 & 12.302 & 163 \\
 & STACK & 2.500 & 12.302 & 206 \\
 & QUEUE & 2.500 & 12.302 & 188 \\
\hline
\multirow{3}{*}{feep.ppm} & Recursivo & 118 & 530 & 9 \\
 & STACK & 118 & 530 & 17 \\
 & QUEUE & 118 & 530 & 17 \\
\hline
\multirow{3}{*}{maze.pbm} & Recursivo & 801 & 4.005 & 39 \\
 & STACK & 801 & 4.005 & 99 \\
 & QUEUE & 801 & 4.005 & 83 \\
\hline
\end{tabular}
\caption{Resultados dos testes de preenchimento de regiões}
\end{table}
Os três algoritmos produziram resultados idênticos em termos de número de acessos à memória (PIXMEM) e pixels preenchidos, confirmando que todas as três implementações estão corretas e apresentam desempenho equivalente em termos de operações de memória.

Relativamente aos tempos de execução observados, a versão recursiva foi consistentemente a mais rápida em todas as imagens testadas. Na imagem $50 \times 50$ branca obteve 163 $\mu$s, seguida pela versão com fila (188 $\mu$s) e pela versão com stack (206 $\mu$s). Este comportamento deve-se ao overhead associado às operações de push/pop em estruturas de dados explícitas comparado com a pilha de sistema. Para a imagem feep.ppm, a versão recursiva foi significativamente mais rápida (9 $\mu$s) que as versões com estruturas explícitas (17 $\mu$s cada). O teste com maze.pbm preencheu 801 pixels do labirinto, com a versão recursiva executando em 39 $\mu$s, fila em 83 $\mu$s e stack em 99 $\mu$s, mantendo o padrão observado de que a versão recursiva é a mais eficiente em termos de tempo de execução.

\section{Segmentação de Imagens}

A função \texttt{ImageSegmentation} identifica e rotula todas as regiões brancas de uma imagem com cores distintas. O algoritmo percorre sequencialmente todos os pixels da imagem. Quando encontra um pixel branco ainda não identificado, gera uma nova cor, aloca-a na LUT, aplica uma função de preenchimento (passada como parâmetro) para colorir a região inteira, e incrementa o contador de regiões. A complexidade temporal é $O(W \times H)$, pois cada pixel é visitado exatamente uma vez no loop principal e pintado no máximo uma vez durante os preenchimentos.

\begin{table}[h]
\begin{tabular}{|l|c|c|c|c|c|}
\hline
\textbf{Teste} & \textbf{Algoritmo} & \textbf{Regiões} & \textbf{PIXMEM} & \textbf{Tempo ($\mu$s)} & \textbf{Cores LUT} \\
\hline
Xadrez $60 \times 60$ & Recursivo & 4 & 11,528 & 88 & 6 \\
Xadrez $60 \times 60$ & STACK & 4 & 11,528 & 140 & 6 \\
Xadrez $60 \times 60$ & QUEUE & 4 & 11,528 & 134 & 6 \\
feep.pbm & Recursivo & 2 & 710 & 10 & 4 \\
feep.pbm & STACK & 2 & 710 & 12 & 4 \\
feep.pbm & QUEUE & 2 & 710 & 13 & 4 \\
maze.pbm & Recursivo & 1 & 5,686 & 48 & 3 \\
maze.pbm & STACK & 1 & 5,686 & 75 & 3 \\
maze.pbm & QUEUE & 1 & 5,686 & 81 & 3 \\
\hline
\end{tabular}
\caption{Resultados da segmentação de imagens com diferentes estratégias}
\end{table}

Os três métodos de preenchimento produziram resultados idênticos em todas as métricas. A LUT final contém as 2 cores originais (branco e preto) mais novas cores geradas para cada região branca identificada. Os valores de PIXMEM confirmam a complexidade $O(W \times H)$ da análise teórica.

A imagem maze.pbm ($41 \times 41$ pixels) representa um labirinto. A segmentação identificou 1 região branca conectada (o caminho do labirinto), com 5,686 acessos PIXMEM em todos os algoritmos, correspondendo à varredura da imagem completa e ao preenchimento da região identificada.

Analisando os tempos de execução da segmentação, observa-se que a versão recursiva foi consistentemente a mais rápida em todos os testes. No xadrez $60 \times 60$ obteve 88 $\mu$s, comparado com 140 $\mu$s (stack) e 134 $\mu$s (fila). Para a imagem feep.pbm, a versão recursiva executou em 10 $\mu$s, enquanto as versões stack e fila demoraram 12 $\mu$s e 13 $\mu$s respetivamente. No labirinto maze.pbm, a versão recursiva foi significativamente mais rápida (48 $\mu$s) que stack (75 $\mu$s) e fila (81 $\mu$s), confirmando que a versão recursiva mantém vantagem de desempenho em todas as escalas de operação.

\section{Conclusões}

Implementaram-se com sucesso todas as funções requisitadas para manipulação de imagens RGB com LUT. As funções de cópia, comparação e rotação apresentam complexidade ótima $O(W \times H)$, sendo impossível fazer melhor considerando que todos os pixels devem ser processados.

A análise detalhada da função \texttt{ImageIsEqual} validou experimentalmente as previsões teóricas, confirmando complexidade $\Theta(W \times H)$ no pior caso. 

Os algoritmos de preenchimento de regiões (flood-fill) foram implementados com sucesso em três variantes. Apesar de apresentarem a mesma complexidade $O(N)$, distinguem-se significativamente em características práticas: a versão recursiva é mais elegante mas com risco de stack overflow; a versão com stack oferece segurança e controlo; a versão com fila oferece uma perspetiva de preenchimento em largura. A seleção do algoritmo apropriado depende do contexto e requisitos.

A função de segmentação integra estes algoritmos para identificar regiões em tempo $O(W \times H)$ e funciona corretamente com qualquer uma das três estratégias de preenchimento, confirmando a generalidade do design implementado.

\end{document}